\section{Atomic Rulings and Robust Declassification}\label{sec:rulings}

\paragraph{One judgment and one dispatch rule.}
Every judgment follows a single canonical form: authority $A$ authorizes engine-rendered call $C$. System vocabulary consists exclusively of mandates, rulings, and log records. A call is dispatched if and only if rulings cover all failed or demanded predicates, each issuing authority's mandate authorizes the exception, and all rulings reference the identical rendered call and are consumed atomically. Multi-party authorization policies (such as two-person control~\cite{saltzer1975protection}) collect multiple rulings for one call. A ruling permits dispatch despite an unsatisfied predicate but never modifies the execution trajectory directly; instead, admitted values update the trajectory Label, and successful tool dispatches record declared side effects in the event log $\mathcal{E}$. No ruling can bypass the explicit agent acceptance required for narrowing (\S\ref{sec:enforcement}).

\paragraph{Atomic plan execution.}
For a plan requiring authorization on a blocked step, a single indivisible operation renders the call, presents provenance (excluding sensitive payload bytes) to the authority, and dispatches upon approval. Plan identifier, ruling, and dispatch event are recorded atomically in the execution log. Acceptance plans are inherently atomic, requiring no external authority round-trip. Approval and dispatch coincide, collapsing traditional multi-stage grant lifecycles (issuing, binding, consuming)~\cite{miller2003paradigm} into an atomic step that prevents call substitution, replay attacks~\cite{needham1978using}, and intermediate interception. In linear logic terms, the capability is consumed simultaneously with its creation~\cite{wadler1990linear}.

\paragraph{Call-scoped release.}
A ruling applies strictly to a single dispatch without altering the trajectory state (the accumulated security label). The ruling and resulting external event record the authorized release, while the security label continues to accurately reflect the underlying classification of the data. This single-command scope isolation ensures an authorized exception or exit cannot persistently relax labels for subsequent attacker-influenced calls~\cite{saltzer1975protection}. Trajectory labels are never persistently widened.

\begin{theorem}[Call-scoped release]\label{thm:release}
In every trajectory: (i)~a call dispatches despite requirement gaps only by atomically consuming, for each gap, a covering ruling over its exact rendered call, and no ruling is consumed by more than one dispatch; (ii)~ruling and acceptance events do not modify any trajectory Label or the committed-effect projection $\mathcal{E}$; hence (iii)~every dispatch with an unmet release-side Label requirement is matched to explicit rulings in the log, and an authorized override widens no later Label and confers no authorization on any other dispatch.
\end{theorem}
Proof in Appendix~\ref{app:proofs}. Corollary~\ref{cor:export} extends the guarantee to child-value export: because labels stay attached to data through merges (Proposition~\ref{prop:preservation}), a later parent release that fails to satisfy the child value's Label requirement requires its own fresh ruling---the formal counterpart of safe data export from child branches, separate from parent-state protection (Theorem~\ref{thm:branch}).

\paragraph{Mandate typology.}
Authority is bound to declarative mandates rather than dynamic runtime objects. Mandates encompass:
(1)~\textbf{Cover mandates}, permitting trust gaps up to a specified rank or audience restrictions up to a defined ceiling;
(2)~\textbf{Named waivers}, explicitly designating exempted \texttt{no\_prior} event types; and
(3)~\textbf{Attends mandates}, specifying attention markers whose requirements a ruling satisfies.
A single ruling may address both an unmet predicate and an attention demand for a call, but cannot bypass narrowing acceptance. Routing uses \texttt{tags}, while security parameters (trust level, audience, side effects) serve as validation criteria; misrouted authorization checks still evaluate against the exact rendered call.

\paragraph{Response-sink rule.}
No end-user ruling can authorize a predicate for a response-sink release, irrespective of mandate privileges, preventing in-band self-approval vulnerabilities~\cite{saltzer1975protection}.

\paragraph{Canonical rulings and robust release.}
Authorities evaluate the engine's canonical representation of a tool call, eliminating semantic evasion vectors where a compromised model requests actions under misleading phrasing. Remedy-set soundness guarantees that every plan generated by the engine is independently sound. Following established declassification models, sanitization regulates *what* information is released, whereas rulings govern *who* permits release~\cite{sabelfeld2009declassification,zdancewic2001robust}. While untrusted input may influence candidate requests formed by the agent, it cannot dictate authorized releases, establishing robust declassification over exact rendered calls.

\paragraph{Sanitizers.}
Registered sanitizers transform values according to authorized state transitions without altering execution trajectories~\cite{caiGhostAgentRedefining2026}. At output boundaries, a derivative is released while the raw output remains confined; at input boundaries, the derived value replaces the raw argument in the canonical tool call. Sanitization adjusts audience scope exclusively and never elevates trust levels without an explicit ruling, type cast, or verified exit transformation (\S\ref{sec:branching}).

\paragraph{User-assistant trust boundary.}
The primary user-assistant conversation trajectory is treated as fully trusted by the engine harness. Because the primary user holds root authority over the agent context, direct conversational replies delivered to the user do not undergo egress declassification filtering. If an interaction involves a recipient who is not trusted with the accumulated security label of the context, that communication must be encapsulated within a tool call (e.g., an outbound messaging API or data export tool). This subjects the payload and action to standard canonical rendering, policy validation, and authorization rulings.

